\documentclass[../00_Main.tex]{subfiles}
\begin{document}
\begin{abstract}
The development and calibration of brain--computer interfaces (\gls{bci}) and electroencephalography (\gls{eeg}) analysis algorithms require controlled and reproducible reference systems. In the absence of a known ground truth in human measurements, phantom heads provide a practical means to study signal transmission, measurement behavior, and signal-processing methods under well-defined conditions. This thesis presents the definition and characterization of a low-cost \gls{eeg} phantom-head system, together with numerical and software tools for controlled signal generation and analysis.

The work introduces a parametric computer-aided design (\gls{cad}) model of an \gls{eeg} phantom head suited for low-cost fabrication. The model defines the head geometry, internal structure, and electrode placement, providing a consistent physical reference that can be used for numerical simulations and experimental configuration.

A saline--gelatine mixture is investigated as a conductive filling material for the phantom head. The electrical properties of this medium are experimentally characterized, demonstrating that its impedance can be tuned through ionic concentration and reproduced across preparations. These results support the feasibility of the mixture as a controllable conductive medium for \gls{eeg} phantom applications.

Based on the defined geometry and material properties, a finite element method (\gls{fem}) pipeline is implemented to model volumetric electrical conduction within the phantom head. The simulations establish quantitative forward and inverse relationships between internal actuation and surface voltage measurements, providing a mathematical basis for analyzing signal propagation and informing the selection of actuation signals for approximating desired surface patterns.

A cross-platform desktop application is developed to generate multi-channel test signals, visualize generated data, and log and export experimental data. Dedicated hardware is used to inject the generated signals from the software platform into the phantom head in a controlled and reproducible manner. The software supports integration with OpenBCI for experimental use.

Together, the \gls{cad} model, conductive material characterization, signal-generation hardware, software platform, and numerical modeling pipeline provide a coherent basis for controlled \gls{eeg} signal simulation and methodological evaluation.
\end{abstract}

\textbf{\textit{Keywords --}} phantom head, EEG simulation, brain--computer interface, finite element method, volumetric conduction, signal generation software

\end{document}